% Input Data Tables for Gujarat Food Grain Distribution Optimization
% Auto-generated by scripts/02_generate_input_tables.py

\begin{table}[htbp]
\centering
\caption{FCI Supply Depots: Monthly Stock and Handling Throughput Capacities (MT)}
\label{tab:supply_capacity}
\resizebox{\textwidth}{!}{
\begin{tabular}{llccccccl}
\toprule
Depot ID & Depot Name & District & Wheat (MT) & Rice (MT) & Total Stock (MT) & Dispatch Cap $C_i$ (MT) & Ratio & Coupling Status \\
\midrule
D1 & FSD Sabarmati (Ahmedabad) & Ahmedabad & 35,000 & 18,000 & 53,000 & 45,000 & 1.18x & Assumed cap ( 8,000 MT below stock) \\
D2 & FSD Chhani (Vadodara) & Vadodara & 22,000 & 12,000 & 34,000 & 28,000 & 1.21x & Assumed cap ( 6,000 MT below stock) \\
D3 & FSD Ghanteshwar (Rajkot) & Rajkot & 25,000 & 10,000 & 35,000 & 30,000 & 1.17x & Assumed cap ( 5,000 MT below stock) \\
D4 & FSD Gandhidham (Kutch) & Kutch & 30,000 & 15,000 & 45,000 & 38,000 & 1.18x & Assumed cap ( 7,000 MT below stock) \\
\midrule
\textbf{TOTAL} & \textbf{Scenario Total} & \textbf{Gujarat} & \textbf{112,000} & \textbf{55,000} & \textbf{167,000} & \textbf{141,000} & \textbf{1.18x} & \textbf{System Cap Margin: +9,000 MT vs Demand} \\
\bottomrule
\end{tabular}
}
\end{table}

\begin{table}[htbp]
\centering
\caption{District Demand Proxies: Assumed Monthly Demand (MT)}
\label{tab:district_demand}
\begin{tabular}{lllcccc}
\toprule
ID & District Warehouse & Facility Location & Wheat (MT) & Rice (MT) & Total (MT) & Share (\%) \\
\midrule
W1 & Ahmedabad District Warehouse & Assumed district demand proxy & 24,000 & 10,000 & 34,000 & 25.76% \\
W2 & Surat District Warehouse & Assumed district demand proxy & 22,000 & 12,000 & 34,000 & 25.76% \\
W3 & Vadodara District Warehouse & Assumed district demand proxy & 14,000 & 7,000 & 21,000 & 15.91% \\
W4 & Rajkot District Warehouse & Assumed district demand proxy & 13,000 & 5,000 & 18,000 & 13.64% \\
W5 & Banaskantha District Warehouse (Palanpur) & Assumed district demand proxy & 11,000 & 5,000 & 16,000 & 12.12% \\
W6 & Kutch District Warehouse (Bhuj) & Assumed district demand proxy & 6,000 & 3,000 & 9,000 & 6.82% \\
\midrule
\textbf{TOTAL} & \textbf{6 Target Districts} & \textbf{State Civil Supplies Grid} & \textbf{90,000} & \textbf{42,000} & \textbf{132,000} & \textbf{100.00%} \\
\bottomrule
\end{tabular}
\end{table}

\begin{table}[htbp]
\centering
\caption{Driving-Profile Distance Matrix $d_{ij}$ (km) and Unit Transport Cost $c_{ij}$ (INR/MT)}
\label{tab:distance_cost_matrix}
\resizebox{\textwidth}{!}{
\begin{tabular}{lcccccc}
\toprule
Depot & W1 (Ahmd) & W2 (Surat) & W3 (Vadodara) & W4 (Rajkot) & W5 (Banas) & W6 (Bhuj) \\
\midrule
D1: Sabarmati (Ahmd) & 6.7 km (\rupee 116.85) & 257.4 km (\rupee 743.55) & 122.3 km (\rupee 405.75) & 220.3 km (\rupee 650.83) & 140.2 km (\rupee 450.50) & 356.9 km (\rupee 992.32) \\
D2: Chhani (Vadodara) & 102.6 km (\rupee 356.42) & 154.9 km (\rupee 487.20) & 7.5 km (\rupee 118.80) & 274.8 km (\rupee 786.88) & 253.3 km (\rupee 733.35) & 450.5 km (\rupee 1226.25) \\
D3: Ghanteshwar (Rajkot) & 222.4 km (\rupee 656.12) & 429.8 km (\rupee 1174.48) & 290.0 km (\rupee 825.08) & 8.5 km (\rupee 121.32) & 326.6 km (\rupee 916.58) & 239.2 km (\rupee 698.10) \\
D4: Gandhidham (Kutch) & 356.7 km (\rupee 991.75) & 564.0 km (\rupee 1510.10) & 424.3 km (\rupee 1160.68) & 199.9 km (\rupee 599.65) & 320.3 km (\rupee 900.70) & 57.6 km (\rupee 244.05) \\
\bottomrule
\end{tabular}
}
\small Note: Values indicate driving-profile route distance in km followed by parenthesized unit cost in INR/MT ($c_{ij} = 100.00 + 2.50 \times d_{ij}$).
\end{table}
